\section*{Introduction}

\begin{center}
\emph{To Oscar Zariski and André Weil.}
\end{center}

\oldpage[I]{5}

This memoir, and the numerous others which are to follow it, are intended to form a treatise on the foundations of Algebraic Geometry.
In principle they presuppose no particular knowledge of this discipline, and it has even turned out that such knowledge, despite its obvious advantages, could sometimes (through the overly exclusive habit of the birational point of view which it implies) be harmful to one who wishes to become familiar with the point of view and the techniques presented here.
On the other hand, we shall assume that the reader has a good knowledge of the following subjects:

\begin{enumerate}
  \item[a)] \emph{Commutative Algebra}, as presented for example in the volumes of the \emph{Éléments} of N.~Bourbaki now in preparation (and, until these volumes appear, in Samuel-Zariski~\cite{I-13} and Samuel~\cite{I-11},~\cite{I-12}).
  \item[b)] \emph{Homological Algebra}, for which we refer to Cartan-Eilenberg~\cite{I-2} (cited as~(M)) and Godement~\cite{I-4} (cited as~(G)), as well as to the recent article of A.~Grothendieck~\cite{I-6} (cited as~(T)).
  \item[c)] \emph{Sheaf Theory}, where our principal references will be~(G) and~(T); this last theory provides the indispensable language for interpreting the essential notions of Commutative Algebra in ``geometric'' terms, and for ``globalizing'' them.
  \item[d)] Finally, it will be useful for the reader to have a certain familiarity with the \emph{functorial language}, which will be used constantly in this Treatise, and for which the reader may consult~(M),~(G) and above all~(T); the principles of this language and the principal results of the general theory of functors will be presented in more detail in a work being prepared by the authors of this Treatise.
\end{enumerate}

\bigskip
\begin{center}*\quad*\quad*\end{center}
\bigskip

This Introduction is not the place to give a more or less summary description of the point of view of ``schemes'' in Algebraic Geometry, nor the long list of reasons which have made its adoption necessary, and in particular the systematic acceptance of nilpotent elements in the local rings of the ``varieties'' that we consider (which necessarily relegates the notion of rational map to the background, in favour of that of regular map or ``morphism'').
The present Treatise aims precisely at developing the language of ``schemes'' systematically, and will, we hope, demonstrate its necessity.
Although it would be easy to do so, we

\oldpage[I]{6}
shall not attempt here either to give an ``intuitive'' introduction to the notions developed in Chapter One.
The reader who would like a preliminary overview of the material of this Treatise may refer to the address given by A.~Grothendieck at the International Congress of Mathematicians in Edinburgh in 1958~\cite{I-7}, and to the exposé~\cite{I-8} by the same author.
The work~\cite{I-14} (cited as~(FAC)) of J.-P.~Serre may also be regarded as an exposition intermediate between the classical point of view and the point of view of schemes in Algebraic Geometry, and as such, reading it may be an excellent preparation for reading our \emph{Éléments}.

\bigskip
\begin{center}*\quad*\quad*\end{center}
\bigskip

For information, we give below the general plan envisaged for this Treatise, which is moreover subject to later modifications, especially as regards the last chapters:

\medskip
\begin{tabular}{rl}
  Chapter One. & --- The language of schemes.\\
  --- \textbf{II}. & --- Elementary global study of some classes of morphisms.\\
  --- \textbf{III}. & --- Cohomology of coherent algebraic sheaves. Applications.\\
  --- \textbf{IV}. & --- Local study of morphisms.\\
  --- \textbf{V}. & --- Elementary procedures for constructing schemes.\\
  --- \textbf{VI}. & --- Descent technique. General method for constructing schemes.\\
  --- \textbf{VII}. & --- Group schemes, principal fibre spaces.\\
  --- \textbf{VIII}. & --- Differential study of fibre spaces.\\
  --- \textbf{IX}. & --- The fundamental group.\\
  --- \textbf{X}. & --- Residues and duality.\\
  --- \textbf{XI}. & --- Intersection theories, Chern classes, Riemann-Roch theorem.\\
  --- \textbf{XII}. & --- Abelian schemes and Picard schemes.\\
  --- \textbf{XIII}. & --- Weil cohomology.\\
\end{tabular}
\medskip

In principle, all the chapters are regarded as open, and supplementary paragraphs may always be added to them later; such paragraphs will appear in separate fascicles, in order to lessen the drawbacks of the mode of publication adopted.
When such a paragraph is planned or in preparation at the time a chapter is published, it will be mentioned in the summary of that chapter, even if, owing to certain orders of urgency, its actual publication should come appreciably later.
For the convenience of the reader, we give in a ``Chapter~0'' various supplements on Commutative Algebra, on Homological Algebra and on Sheaf Theory, used in the course of the chapters of this Treatise, which are more or less well known, but for which it has not been possible to give convenient references.
The reader is advised to refer to Chapter~0 only in the course of reading the Treatise proper, and to the extent that the results to which we refer are not

\oldpage[I]{7}
sufficiently familiar to the reader.
We believe, moreover, that in this way the reading of this Treatise may be a good method for the beginner to become familiar with Commutative Algebra and Homological Algebra, the study of which, when it is not accompanied by tangible applications, is judged tedious, even depressing, by a fairly large number of people.

\bigskip
\begin{center}*\quad*\quad*\end{center}
\bigskip

It is beyond our competence to give in this Introduction a historical survey, even a summary one, of the notions and results presented.
The text will contain only the references judged particularly useful for its understanding, and we shall indicate the origin only of the most important results.
Formally at least, the subjects treated in our work are fairly new, which will explain the scarcity of references to the Fathers of Algebraic Geometry of the nineteenth century and of the beginning of the twentieth century, whose works we know only by hearsay.
It is appropriate, however, to say a few words here about the works which have most directly influenced the authors and contributed to the development of the point of view of schemes.
First and foremost we must cite the fundamental work (FAC) of J.-P.~Serre, which has served as an introduction to Algebraic Geometry for more than one young adept (among them one of the authors of the present Treatise), put off by the aridity of A.~Weil's classic \emph{Foundations}~\cite{I-18}.
That work was the first to show that the ``Zariski topology'' of an ``abstract'' algebraic variety lends itself perfectly to certain techniques of Algebraic Topology, and in particular to the construction of a cohomology theory.
Moreover, the definition of an algebraic variety given there is the one which lends itself most naturally to the extension of this notion that we develop here\footnote{As J.-P.~Serre has pointed out to us, it should be noted that the idea of defining the structure of a variety through a sheaf of rings goes back to H.~Cartan, who made it the starting point of his theory of analytic spaces.
Naturally, in ``Analytic Geometry'' just as in Algebraic Geometry, nilpotent elements ought to be given full rights of citizenship in the local rings of analytic spaces.
This extension of the definition of H.~Cartan and J.-P.~Serre has recently been taken up by H.~Grauert~\cite{I-5}, and there is reason to hope that a systematic exposition of Analytic Geometry in this general framework will soon see the light of day.
It is moreover evident that the notions and techniques developed in this Treatise retain a meaning in Analytic Geometry, although one must expect more considerable technical difficulties in the latter theory.
One may foresee that Algebraic Geometry, through the simplicity of its methods, will be able to serve as a kind of formal model for future developments in the theory of analytic spaces.}.
Serre had moreover observed himself that the cohomology theory of affine algebraic varieties carries over without difficulty when the affine algebras over a field are replaced by arbitrary commutative rings.
In essence, therefore, Chapters I and II of this Treatise and the first two paragraphs of Chapter III can be viewed as easy transpositions, to this enlarged framework, of the principal results of (FAC) and of a later article by the same author~\cite{I-15}.
We have also profited greatly from the \emph{Séminaire de Géométrie algébrique} of C.~Chevalley~\cite{I-1}; in particular, the systematic use of the ``constructible sets'' introduced by him has proved very useful in the theory of schemes (cf.\ Chapter~IV).
From him, too, we have taken the study of morphisms from the point

\oldpage[I]{8}
of view of dimension (Chapter~IV), which carries over without notable change to the framework of schemes.
It should be noted, furthermore, that the notion of ``schemes of local rings'', introduced by Chevalley, lends itself naturally to an extension of Algebraic Geometry (not having, however, all the flexibility and generality that we intend to give it here); for the relations between this notion and our theory, see Chapter One, \S~8.
Such an extension has been developed by M.~Nagata in a series of memoirs~\cite{I-9} containing numerous special results concerning Algebraic Geometry over Dedekind rings\footnote{Among the works which come close to our point of view in Algebraic Geometry, let us mention the important work of E.~Kähler~\cite{I-22}, and a recent Note of Chow and Igusa~\cite{I-3}, who take up, within the framework of the Nagata-Chevalley theory, certain results of (FAC) and also give a Künneth formula.}.

\bigskip
\begin{center}*\quad*\quad*\end{center}
\bigskip

Finally, it goes without saying that a book on Algebraic Geometry, and above all a book dealing with foundations, is necessarily influenced, if only through intermediaries, by mathematicians such as O.~Zariski and A.~Weil.
In particular, the \emph{Theory of holomorphic functions} of Zariski~\cite{I-20}, suitably made more flexible by means of cohomological methods and completed by an existence theorem (Chapter~III, \S\S~4 and 5), is (together with the descent technique presented in Chapter~VI) one of the principal tools used in this Treatise, and seems to us one of the most powerful available in Algebraic Geometry.

The general technique into which it fits can be sketched as follows (a typical example of it will be provided in Chapter~IX, in the study of the fundamental group).
One has a proper morphism (Chapter~II) $f:X\to Y$ from one algebraic variety to another (more generally, from one scheme to another), which one wishes to study in a neighbourhood of a point $y\in Y$, with a view to solving a problem P relative to a neighbourhood of $y$.
One proceeds in successive steps:
\begin{enumerate}
  \item[$1^\circ$] One may suppose $Y$ affine, so that $X$ becomes a scheme defined over the affine ring $A$ of $Y$, and one may even replace $A$ by the local ring of $y$.
  This reduction is always easy in practice (Chapter~V) and brings us back to the case where $A$ is a \emph{local} ring.
  \item[$2^\circ$] One studies the problem under consideration when $A$ is an \emph{Artinian} local ring.
  In order for it to retain a genuine meaning when $A$ is not assumed to be integral, it is sometimes necessary to reformulate the problem P, and it turns out that one often thereby obtains a better understanding of the problem, which at this stage is of an ``infinitesimal'' nature.
  \item[$3^\circ$] The theory of formal schemes (Chapter~III, \S\S~3, 4 and 5) makes it possible to pass from the case of an Artinian ring to the case of a \emph{complete local ring}.
  \item[$4^\circ$] Finally, if $A$ is an arbitrary local ring, the consideration of ``multivalued sections'' over suitable schemes over $X$, approximating a given ``formal'' section (Chapter~IV), will often make it possible to pass from a result known for the scheme

  \oldpage[I]{9}
  deduced from $X$ by extension of scalars to the completion of $A$, to an analogous result for a fairly simple (for example unramified) finite extension of $A$.
\end{enumerate}

This sketch shows the importance of the systematic study of schemes defined over an Artinian ring $A$.
Serre's point of view in his formulation of local class field theory, and recent work of Greenberg, seem to suggest that such a study could be undertaken by functorially attaching to such a scheme $X$ a scheme $X'$ over the residue field $k$ of $A$ (assumed perfect), of dimension equal (in favourable cases) to $n\dim X$, where $n$ is the length of $A$.

As for the influence of A.~Weil, suffice it to say that it is the necessity of developing the tools needed to formulate, with all the desired generality, the definition of ``Weil cohomology'', and to undertake the proof\footnote{To avoid any misunderstanding, let us make clear that this task has barely been undertaken at the time this Introduction is being written, and has therefore not yet led to the proof of the Weil conjectures.} of all the formal properties needed to establish his celebrated conjectures in Diophantine Geometry~\cite{I-19}, that has been one of the principal motivations for the writing of the present Treatise, on the same footing as the desire to find the natural framework for the usual notions and methods of Algebraic Geometry, and to give the authors the opportunity of understanding the said notions and techniques.

\bigskip
\begin{center}*\quad*\quad*\end{center}
\bigskip

To conclude, we think it useful to warn readers that, just like the authors themselves, they will no doubt have some difficulty before becoming accustomed to the language of schemes, and before convincing themselves that the usual constructions suggested by geometric intuition can be transcribed, in essentially only one reasonable way, into this language.
As in many parts of modern Mathematics, the initial intuition moves further and further away, in appearance, from the language suited to expressing it with all the desired precision and generality.
In the present case, the psychological difficulty stems from the necessity of transferring to the objects of a category already rather different from the category of sets (namely the category of preschemes, or the category of preschemes over a given prescheme) notions which are familiar for sets: Cartesian products, laws of group, of ring, of module, fibre bundles, principal homogeneous fibre bundles, etc.
It will no doubt be difficult for the mathematician, in the future, to shirk this new effort of abstraction, perhaps rather slight, all things considered, in comparison with the one made by our fathers in familiarizing themselves with Set Theory.

\bigskip
\begin{center}*\quad*\quad*\end{center}
\bigskip

References will be given according to the decimal system; for example, in III, 4.9.3, the numeral III indicates the chapter, the numeral 4 the paragraph, and the numeral 9 the section of the paragraph.
Within the same chapter, mention of the chapter will be omitted.
